\documentclass[11pt]{article}

\usepackage[T1]{fontenc}
\usepackage[utf8]{inputenc}
\usepackage{amsmath,amssymb,amsthm,mathtools}
\usepackage{microtype}
\usepackage{hyperref}
\hypersetup{hidelinks}
\usepackage{geometry}
\geometry{margin=1in}

\newtheorem{theorem}{Theorem}[section]
\newtheorem{proposition}[theorem]{Proposition}
\newtheorem{lemma}[theorem]{Lemma}
\newtheorem{corollary}[theorem]{Corollary}
\newtheorem{remark}[theorem]{Remark}
\newtheorem{definition}[theorem]{Definition}

\newcommand{\C}{\mathbb{C}}
\newcommand{\op}{\mathrm{op}}
\newcommand{\cmax}{C_{\max}}
\newcommand{\dd}{\delta}
\newcommand{\etaq}{\eta}
\newcommand{\tauq}{\tau}
\newcommand{\PhiRho}{\Phi_{\rho}}

\title{Sharp Stability and Two-Cell Rigidity for\\
Finite Projector--Unitary Commutator Profiles}

\author{Bertrand Dalimier\\
Independent Researcher}

\date{}

\begin{document}

\maketitle

\begin{abstract}
Let $(P_i)_{i\in I}$ be a finite orthogonal resolution of the identity on a
finite-dimensional complex Hilbert space and let $U$ be unitary.  Write
$a_i=\|[P_i,U]\|_{\op}$,
$\dd^2=\sum_i a_i^2$, and
\[
  \cmax=\max_{S\subseteq I}
  \left\|\left[\sum_{i\in S}P_i,U\right]\right\|_{\op}.
\]
We prove the sharp finite-family envelope
$\cmax^2\le\dd^2/2$.  For $\dd>0$, equality holds if and only if exactly two
cell commutators are nonzero.  We then prove a dimension-free inverse
near-equality theorem: if
\[
  \cmax^2\ge(1-\etaq)\frac{\dd^2}{2},
  \qquad
  0\le\etaq\le\frac15,
  \qquad
  \dd\ge\rho>0,
\]
then two cells carry all but a relative fraction
\[
  2\etaq+\frac{36\etaq}{(1-3\etaq)\rho^2}
\]
of the global quadratic defect.  For $\etaq\le1/6$ this is at most
$\etaq(2+72/\rho^2)$.  An explicit three-cell family proves that the
inverse-square defect-scale dependence is necessary up to multiplicative
constants, with
\[
  \lim_{m\to\infty}\left[
    \lim_{n\to\infty}
    \frac{\tauq_{n,m}}{\etaq_{n,m}}\dd_{n,m}^2
  \right]=4.
\]
The principal theorems and sharpness construction are machine-checked in
Lean 4.  They form the downstream layer of an audited end-to-end
finite-dimensional quantum-foundations library that also formalizes the
Busch and Gleason representation theorems, Naimark dilation, Wigner's theorem,
Uhlhorn-type uniqueness, Riedel's branch decomposition, and Kent's
contrary-inference construction.
\end{abstract}

\noindent\textbf{Keywords.}
projector--unitary commutators; near-equality stability; rigidity;
Lean 4 formal verification; Gleason theorem; Busch theorem;
Kent contrary inference

\noindent\textbf{2020 Mathematics Subject Classification.}
15A60 (primary); 15A45, 47A63 (secondary).

\section{Introduction}

Let a finite-dimensional Hilbert space be decomposed by a finite orthogonal
resolution of the identity $(P_i)_{i\in I}$.  A unitary $U$ preserves the
$i$th summand exactly when $[P_i,U]=0$, so the commutator norms provide a
natural finite defect profile for the failure of $U$ to respect the
decomposition.  The individual operator norms
\[
  a_i:=\|[P_i,U]\|_{\op}
\]
therefore define a nonnegative finite defect profile.  The quadratic aggregate
\[
  \dd^2:=\sum_{i\in I} a_i^2
\]
is a natural dimension-free global measure of failure of simultaneous
commutation.

A second object is obtained by grouping sectors into a cut.  For
$S\subseteq I$, let
\[
  P_S:=\sum_{i\in S}P_i,\qquad
  C_S:=\|[P_S,U]\|_{\op}.
\]
Since $I$ is finite, the cut envelope
\[
  \cmax:=\max_{S\subseteq I}C_S
\]
is attained.  The basic inequality studied here is
\begin{equation}\label{eq:envelope}
  \cmax^2\le \frac{\dd^2}{2}.
\end{equation}
It is a finite operator-theoretic inequality: no probability assignment,
dynamical assumption, open-system structure, or asymptotic Hilbert-space
limit is involved.

The purpose of this paper is to analyze the near-equality regime of
\eqref{eq:envelope}.  There are three levels.

First, we identify the exact equality cases.  For positive global defect,
saturation of \eqref{eq:envelope} is equivalent to the statement that exactly
two cells are active: all other cellwise commutator defects vanish.

Second, we prove a quantitative inverse theorem.  If the cut envelope nearly
saturates \eqref{eq:envelope}, then almost all of the quadratic defect budget
must be carried by two cells.  The estimate is independent of the number of
cells.  At a fixed positive defect scale $\rho$, the relative two-cell tail is
$O(\etaq/\rho^2)$ as the saturation deficit $\etaq$ tends to zero.

Third, we construct an explicit three-cell family showing that the
$1/\rho^2$ scale is not a proof artifact.  The family has closed formulas for
the three cell budgets, the global defect, the maximal cut, the
near-saturation parameter, and the two-cell tail.  Its asymptotics give the
constant
\[
  \lim_{m\to\infty}\left[
    \lim_{n\to\infty}\frac{\tauq_{n,m}}{\etaq_{n,m}}\dd_{n,m}^2
  \right]=4.
\]
Thus a stability estimate uniform down to zero defect cannot in general have
a tail bound of order $O(\etaq)$ with a scale-independent constant.

The arguments separate into a finite real concentration layer and an
operator-theoretic layer.  The principal results are machine checked in a
companion Lean 4 development.  This development is the downstream component
of a pinned four-repository quantum-foundations stack whose audited upstream
coverage includes Busch, Gleason, Naimark, Wigner, Uhlhorn-type uniqueness,
Riedel's branch decomposition, and Kent's contrary-inference construction.
These results provide formalization context and provenance; they are not
premises of the cut-envelope theorem unless explicitly stated.

% TODO before submission:
% 1. Replace repository commit by archival release/DOI.
% 2. Apply final LAA/Elsevier template if required.
% 3. Prepare cover letter and submission metadata.

\section{Finite orthogonal-resolution setup}

Let $H$ be a finite-dimensional complex Hilbert space.  Let
$I$ be a finite index set and let $(P_i)_{i\in I}$ be orthogonal projections
satisfying
\begin{equation}\label{eq:decomp}
  P_iP_j=0\quad(i\neq j),
  \qquad
  \sum_{i\in I}P_i=I_H.
\end{equation}
We call the family a finite orthogonal resolution of the identity.  Let
$U:H\to H$ be unitary.  In the broader quantum-foundations development from
which this problem arose, the same objects are also used as finite projective
record decompositions; that interpretation is not needed for the results
below.

For each cell $i$, define
\begin{equation}\label{eq:ai}
  a_i:=\|[P_i,U]\|_{\op},
  \qquad
  p_i:=a_i^2,
\end{equation}
and define the global quadratic commutator defect
\begin{equation}\label{eq:delta}
  \dd:=\left(\sum_{i\in I}p_i\right)^{1/2}.
\end{equation}
For a cut $S\subseteq I$, set
\begin{equation}\label{eq:cut}
  P_S:=\sum_{i\in S}P_i,
  \qquad
  C_S:=\|[P_S,U]\|_{\op},
  \qquad
  \cmax:=\max_{S\subseteq I}C_S.
\end{equation}

The quantity $\dd$ is the $\ell^2$ norm of the cellwise operator-norm profile.
The quantity $\cmax$ instead probes the largest coherent defect visible after
aggregating an arbitrary subset of the cells.

For later use, suppose $\dd>0$.  Then the decomposition has at least two
cells: a one-cell decomposition would have $P_1=I_H$ and hence zero
commutator defect.  We may therefore define the optimal two-cell tail
\begin{equation}\label{eq:tau}
  \tauq:=
  \min_{\substack{i,j\in I\\i\neq j}}
  \frac{\dd^2-p_i-p_j}{\dd^2}.
\end{equation}
Equivalently, $\tauq$ is the fraction of the total quadratic budget left after
retaining the best pair of distinct cells.

\begin{remark}
The formal theorem is stated without division by $\dd^2$: it asserts the
existence of distinct $i,j$ such that
\[
  \dd^2-p_i-p_j\le \varepsilon\,\dd^2.
\]
This formulation remains meaningful at zero defect.  In the positive-defect
regime it is equivalent to $\tauq\le\varepsilon$.
\end{remark}

\section{The finite cut envelope}

The key observation is that a cut commutator is controlled simultaneously by
the defect budget on the cut and by the defect budget on its complement.

Relative to the orthogonal splitting
\[
  H=P_SH\oplus (I-P_S)H,
\]
the commutator $[P_S,U]$ is off diagonal.  Its two cross blocks are
\[
  P_SU(I-P_S),
  \qquad
  (I-P_S)UP_S.
\]
For a unitary $U$ the operator norm of the cut commutator is the maximum of
the norms of these two cross blocks.  The incoming block is controlled by the
cellwise budgets on $S$, while the outgoing block is controlled by the
budgets on $S^c$.  Consequently,
\begin{equation}\label{eq:cut-min}
  C_S^2
  \le
  \min\left\{
    \sum_{i\in S}p_i,\,
    \sum_{i\notin S}p_i
  \right\}.
\end{equation}

Since the two quantities inside the minimum sum to $\dd^2$, we obtain the
universal envelope.

\begin{theorem}[Cut-envelope inequality]\label{thm:envelope}
For every finite orthogonal resolution and every unitary $U$,
\[
  \boxed{\cmax^2\le\frac{\dd^2}{2}}.
\]
\end{theorem}

\begin{proof}
For every $S\subseteq I$, \eqref{eq:cut-min} gives
\[
  C_S^2
  \le
  \min\{A_S,\dd^2-A_S\},
  \qquad
  A_S:=\sum_{i\in S}p_i.
\]
The elementary inequality
$\min\{x,T-x\}\le T/2$ for $0\le x\le T$ yields
$C_S^2\le\dd^2/2$.  Taking the maximum over the finite set of cuts proves the
claim.
\end{proof}

The coefficient $1/2$ is optimal.  More importantly for the present paper,
its equality cases are rigid.

\section{Quantitative stability of near-saturators}

We now prove the inverse theorem behind Theorem~\ref{thm:rigidity}.  For a
cut $S\subseteq I$, write
\[
  A_S:=\sum_{i\in S}p_i,
  \qquad
  C_S:=\|[P_S,U]\|_{\op}.
\]

\subsection{A one-sided concentration lemma}

The numerical constant in the final stability theorem enters through the
following dimension-free lemma.

\begin{lemma}[One-sided concentration]\label{lem:onesided}
Let $S$ be a cut with $C_S>0$, and put
\[
  e_S:=A_S-C_S^2.
\]
If $e_S\le C_S^2/2$, then there exists $i\in S$ such that
\begin{equation}\label{eq:onesided}
  A_S-p_i
  \le
  2e_S+\frac{18e_S}{C_S^2-e_S}.
\end{equation}
\end{lemma}

\begin{proof}
Relative to $P_SH\oplus P_{S^c}H$, the cut commutator has incoming and
outgoing cross blocks.  We give the argument when the incoming block realizes
$C_S$; the outgoing case follows by the same argument after exchanging the
two blocks.

Because the space is finite dimensional, choose a unit vector
$v\in P_{S^c}H$ on which the incoming block attains its norm.  For $c\in S$
set
\[
  z_c:=U^*P_cUv,\qquad
  r_c:=\|z_c\|^2,\qquad
  \alpha_c:=p_c-r_c.
\]
The vectors $z_c$ are mutually orthogonal and the witness identities give
\begin{equation}\label{eq:witness-defect}
  \alpha_c\ge0,\qquad
  \sum_{c\in S}\alpha_c=e_S.
\end{equation}
Call a cell good when $\alpha_c\le p_c/2$.  The total $p$-mass of the bad
cells is at most
\begin{equation}\label{eq:badmass}
  \sum_{\text{bad }c}p_c\le2e_S.
\end{equation}
Discarding good cells with $p_c=0$ does not change the good-cell mass.  On
the remaining active good cells one has $r_c>0$.  Define
\[
  \xi_c:=\frac{z_c}{\sqrt{r_c}}
\]
and a normalized comparison family $\zeta_c$ obtained by replacing the
residual component of $\xi_c$ by $\sqrt{p_c}\,v$.  The witness geometry gives
three facts:
\begin{align}
  &(\xi_c)_c\ \text{is orthonormal}, \label{eq:xi-orth}\\
  &p_cp_d=|\langle\zeta_c,\zeta_d\rangle|^2
      \quad(c\ne d), \label{eq:zeta-gram}\\
  &\sum_c\|\zeta_c-\xi_c\|^2\le2e_S. \label{eq:gram-close}
\end{align}

For completeness, if $\xi$ is any orthonormal family and $\zeta$ satisfies
\eqref{eq:zeta-gram}, with
$q:=\sum_c\|\zeta_c-\xi_c\|^2$, expansion of each off-diagonal inner product
around the orthonormal reference family and Bessel's inequality give
\begin{equation}\label{eq:gram-spread-general}
  \sum_c\sum_{d\ne c}p_cp_d
  \le 6q+3q^2.
\end{equation}
Here $C_S\le1$, hence $e_S\le1/2$.  By \eqref{eq:gram-close},
$q\le2e_S$, so
\begin{equation}\label{eq:gram18}
  \sum_c\sum_{d\ne c}p_cp_d\le18e_S.
\end{equation}
For a finite nonnegative profile,
\[
  \sum_c\sum_{d\ne c}p_cp_d
  =
  \left(\sum_cp_c\right)^2-\sum_cp_c^2.
\]
Choosing a maximal entry therefore shows that if the total active-good mass
is at least $L>0$, one active good cell leaves active-good tail at most
$18e_S/L$.

The witness identities and \eqref{eq:witness-defect} give the lower bound
\[
  L=C_S^2-e_S
\]
for that active-good mass.  Adding the bad-cell contribution
\eqref{eq:badmass} yields \eqref{eq:onesided}.
\end{proof}

\subsection{Bilateral concentration across a cut}

Let
\[
  A:=\dd^2,\qquad C:=C_S,\qquad
  A_T:=A_{S^c},
\]
and define
\[
  e_S:=A_S-C^2,\qquad
  e_T:=A_T-C^2,\qquad
  E:=A-2C^2.
\]
The complementary budget identity $A_S+A_T=A$ gives
\begin{equation}\label{eq:e-split}
  e_S+e_T=E.
\end{equation}

\begin{proposition}[Bilateral cut concentration]\label{prop:bilateral}
Assume $C>0$ and $E\le C^2/2$.  Then there exist
$i\in S$ and $j\in S^c$ such that
\begin{equation}\label{eq:bilateral}
  A-p_i-p_j
  \le
  2E+\frac{18E}{C^2-E}.
\end{equation}
\end{proposition}

\begin{proof}
The cut estimate gives $C^2\le A_S$ and $C^2\le A_T$, hence
$e_S,e_T\ge0$.  By \eqref{eq:e-split}, each is at most $E$ and therefore at
most $C^2/2$.  Lemma~\ref{lem:onesided}, applied to $S$ and $S^c$, gives
$i\in S$ and $j\in S^c$ such that
\[
  A_S-p_i
  \le
  2e_S+\frac{18e_S}{C^2-e_S},
\]
\[
  A_T-p_j
  \le
  2e_T+\frac{18e_T}{C^2-e_T}.
\]
Set $d:=C^2-E>0$.  Since
$d\le C^2-e_S$ and $d\le C^2-e_T$,
\[
  \frac{e_S}{C^2-e_S}
  +
  \frac{e_T}{C^2-e_T}
  \le
  \frac{e_S+e_T}{d}
  =
  \frac Ed.
\]
Adding the two one-sided inequalities proves
\eqref{eq:bilateral}.  Since $i\in S$ and $j\in S^c$, the two cells are
distinct.
\end{proof}

\subsection{The scale-sensitive modulus}

Let $0\le\etaq\le1/5$ and assume
\begin{equation}\label{eq:near}
  \cmax^2
  \ge
  (1-\etaq)\frac{\dd^2}{2}.
\end{equation}
We also impose the positive scale condition
\begin{equation}\label{eq:scale}
  \dd\ge\rho>0.
\end{equation}
Define
\begin{equation}\label{eq:modulus}
  \PhiRho(\etaq):=
  2\etaq+
  \frac{36\etaq}{(1-3\etaq)\rho^2}.
\end{equation}

\begin{theorem}[Quantitative two-cell stability]\label{thm:stability}
Suppose \eqref{eq:near} and \eqref{eq:scale} hold with
$0\le\etaq\le1/5$.  Then there exist distinct cells $i,j$ such that
\begin{equation}\label{eq:stability-conclusion}
  \boxed{
  \dd^2-p_i-p_j
  \le
  \PhiRho(\etaq)\,\dd^2.
  }
\end{equation}
Equivalently,
\[
  \tauq\le
  2\etaq+
  \frac{36\etaq}{(1-3\etaq)\rho^2}.
\]
\end{theorem}

\begin{proof}
Choose a cut $S$ attaining $\cmax$ and put
\[
  A:=\dd^2,\qquad C:=\cmax,\qquad E:=A-2C^2.
\]
Theorem~\ref{thm:envelope} gives $E\ge0$, while
\eqref{eq:near} gives
\begin{equation}\label{eq:Eeta}
  E\le\etaq A.
\end{equation}
Since $\etaq\le1/5$ and $A=2C^2+E$,
\[
  E\le\frac A5
  \quad\Longrightarrow\quad
  E\le\frac{C^2}{2}.
\]
Moreover $A>0$, and the same inequalities imply $C>0$.  Proposition
\ref{prop:bilateral} therefore gives distinct $i,j$ with
\[
  A-p_i-p_j
  \le
  2E+\frac{18E}{d},
  \qquad
  d:=C^2-E.
\]
Because $C^2=(A-E)/2$,
\begin{equation}\label{eq:didentity}
  d=\frac{A-3E}{2}.
\end{equation}
Using \eqref{eq:Eeta},
\[
  d\ge\frac{(1-3\etaq)A}{2}>0
\]
and hence
\[
  \frac Ed
  \le
  \frac{2\etaq}{1-3\etaq}.
\]
Thus
\begin{equation}\label{eq:absolute-tail}
  A-p_i-p_j
  \le
  2\etaq A+
  \frac{36\etaq}{1-3\etaq}.
\end{equation}
Finally $A=\dd^2\ge\rho^2$, so
\[
  \frac{36\etaq}{1-3\etaq}
  \le
  A\,
  \frac{36\etaq}{(1-3\etaq)\rho^2}.
\]
Substitution into \eqref{eq:absolute-tail} gives
\eqref{eq:stability-conclusion}.
\end{proof}

At every fixed positive scale,
\[
  \PhiRho(\etaq)\longrightarrow0
  \qquad(\etaq\downarrow0).
\]

\begin{corollary}[Transparent linear form]\label{cor:linear}
If in addition $\etaq\le1/6$, then
\begin{equation}\label{eq:linear}
  \boxed{
    \tauq\le
    \etaq\left(2+\frac{72}{\rho^2}\right).
  }
\end{equation}
\end{corollary}

\begin{proof}
For $\etaq\le1/6$ one has
$1-3\etaq\ge1/2$, so
\[
  \frac{36\etaq}{(1-3\etaq)\rho^2}
  \le
  \frac{72\etaq}{\rho^2}.
\]
Apply Theorem~\ref{thm:stability}.
\end{proof}


\section{Exact rigidity at saturation}

Call a cell \emph{active} when $p_i>0$.

\begin{theorem}[Exact equality classification]\label{thm:rigidity}
Assume $\dd>0$.  Then
\[
  \boxed{
    \cmax^2=\frac{\dd^2}{2}
    \quad\Longleftrightarrow\quad
    \text{exactly two cells are active}.
  }
\]
More explicitly, equality holds if and only if there exist distinct
$i,j\in I$ such that
\[
  p_i>0,\qquad p_j>0,\qquad
  p_k=0\quad\text{for every }k\notin\{i,j\}.
\]
\end{theorem}

\begin{proof}
Set $A=\dd^2$.  First assume saturation,
$\cmax^2=A/2$.  Apply Theorem~\ref{thm:stability} with
$\etaq=0$ and $\rho=\dd>0$.  Since $\Phi_{\dd}(0)=0$, there are distinct
$i,j$ such that
\[
  A-p_i-p_j\le0.
\]
On the other hand,
\[
  A-p_i-p_j=\sum_{k\notin\{i,j\}}p_k\ge0,
\]
so every $p_k$ with $k\notin\{i,j\}$ vanishes.

It remains to show that the two selected cells are genuinely active.  A
singleton is an admissible cut, and therefore
\[
  p_i=\|[P_i,U]\|_{\op}^2\le\cmax^2=\frac A2,
  \qquad
  p_j\le\frac A2.
\]
Since the remaining cells carry zero budget, $p_i+p_j=A$.  Hence
$p_i=p_j=A/2>0$.

Conversely, suppose exactly two cells $i\ne j$ are active.  Then
$A=p_i+p_j$, so
\[
  \max\{p_i,p_j\}\ge\frac A2.
\]
Again using singleton cuts,
\[
  \cmax^2\ge\max\{p_i,p_j\}\ge\frac A2.
\]
The cut-envelope inequality, Theorem~\ref{thm:envelope}, gives the reverse
bound $\cmax^2\le A/2$, and equality follows.
\end{proof}

The hypothesis $\dd>0$ only removes the degenerate configuration in which
every commutator vanishes.


\section{An explicit three-cell sharpness family}\label{sec:sharp}

We construct a family in dimension three whose near-saturation deficit tends
to zero while the ratio $\tauq/\etaq$ diverges at the inverse-square defect
scale.

Let $(p_n,q_n)$ be the positive integer recurrence
\[
  (p_0,q_0)=(1,1),
\]
\[
  p_{n+1}=3p_n+4q_n,
  \qquad
  q_{n+1}=2p_n+3q_n.
\]
It preserves
\begin{equation}\label{eq:pell}
  p_n^2+1=2q_n^2.
\end{equation}
Define
\[
  a_n=\frac{p_n+1}{2q_n},
  \qquad
  b_n=\frac{p_n-1}{2q_n},
\]
so that $a_n^2+b_n^2=1$.  Rotate these amplitudes by $45^\circ$:
\begin{equation}\label{eq:AB}
  A_n=\frac{a_n+b_n}{\sqrt2},
  \qquad
  B_n=\frac{a_n-b_n}{\sqrt2}.
\end{equation}
Then
\[
  A_n^2+B_n^2=1,\qquad
  B_n=\frac{1}{\sqrt2\,q_n},
\]
hence
\begin{equation}\label{eq:AB-limits}
  B_n^2\to0,\qquad A_n^2\to1.
\end{equation}

For a second index $m$, let $c_m,s_m$ be the explicit rational unit-circle
parameters
\[
  c_m^2+s_m^2=1,\qquad s_m>0,
\]
coming from the same Pell family, and set
\begin{equation}\label{eq:d}
  d_m:=1-c_m.
\end{equation}
Thus
\begin{equation}\label{eq:s-d}
  s_m^2=2d_m-d_m^2,
  \qquad
  d_m\to0.
\end{equation}

The three-dimensional unitary is obtained by conjugating the planar rotation
with a basis change whose first direction is
$(A_n,B_n,0)$.  For the three rank-one cells, the exact squared commutator
budgets are
\begin{align}
  p_0 &=
  2d_mA_n^2-d_m^2A_n^4, \label{eq:p0}\\
  p_1 &=
  2d_mB_n^2-d_m^2B_n^4, \label{eq:p1}\\
  p_2 &= s_m^2. \label{eq:p2}
\end{align}
These identities follow from the rank-one formula
\begin{equation}\label{eq:rankone}
  \boxed{
  \|[P_e,U]\|_{\op}^2
  =
  1-\bigl|\langle e,Ue\rangle\bigr|^2
  }
\end{equation}
for a unit vector $e$ and the rank-one projector $P_e$.

Summing \eqref{eq:p0}--\eqref{eq:p2} gives
\begin{equation}\label{eq:global-sharp}
  \boxed{
  \dd_{n,m}^2
  =
  2s_m^2+
  2d_m^2A_n^2B_n^2.
  }
\end{equation}
The maximal cut is exact:
\begin{equation}\label{eq:cmax-sharp}
  \boxed{
  \cmax(n,m)=s_m.
  }
\end{equation}

Define the relative saturation deficit
\[
  \etaq_{n,m}
  :=
  \frac{\dd_{n,m}^2-2\cmax(n,m)^2}{\dd_{n,m}^2}.
\]
Then
\begin{equation}\label{eq:eta-sharp}
  \boxed{
  \etaq_{n,m}
  =
  \frac{
    2d_m^2A_n^2B_n^2
  }{
    2s_m^2+2d_m^2A_n^2B_n^2
  }.
  }
\end{equation}
The cell-$1$ budget is the smallest of the three:
\begin{equation}\label{eq:cell1-min}
  p_1\le p_0,
  \qquad
  p_1\le p_2.
\end{equation}
Indeed, $B_n^2\le A_n^2$ and the exact budget formulas imply the first
inequality, while the maximal-cut estimate for the family gives the second.
Consequently, retaining cells $0$ and $2$ is the optimal two-cell choice in
the sense of \eqref{eq:tau}.  Thus the general optimal tail is exactly
\begin{equation}\label{eq:tau-sharp}
  \boxed{
  \tauq_{n,m}
  =
  \frac{
    2d_mB_n^2-d_m^2B_n^4
  }{
    2s_m^2+2d_m^2A_n^2B_n^2
  }.
  }
\end{equation}
This identification is important: the sharpness calculation below concerns
the same optimal two-cell tail that appears in
Theorem~\ref{thm:stability}, not merely a particular candidate pair.
The exact quotient simplifies to
\begin{equation}\label{eq:ratio-sharp}
  \boxed{
  \frac{\tauq_{n,m}}{\etaq_{n,m}}
  =
  \frac{2-d_mB_n^2}{2d_mA_n^2}.
  }
\end{equation}

\begin{theorem}[Scale sharpness]\label{thm:scale-sharp}
For fixed $m$,
\begin{equation}\label{eq:fixed-m-ratio}
  \lim_{n\to\infty}
  \frac{\tauq_{n,m}}{\etaq_{n,m}}
  =
  \frac1{d_m},
\end{equation}
and
\begin{equation}\label{eq:fixed-m-delta}
  \lim_{n\to\infty}\dd_{n,m}^2=2s_m^2.
\end{equation}
Consequently,
\begin{equation}\label{eq:iterated4}
  \boxed{
  \lim_{m\to\infty}
  \left[
    \lim_{n\to\infty}
    \frac{\tauq_{n,m}}{\etaq_{n,m}}
    \dd_{n,m}^2
  \right]
  =4.
  }
\end{equation}
\end{theorem}

\begin{proof}
Equations \eqref{eq:fixed-m-ratio} and \eqref{eq:fixed-m-delta} follow from
\eqref{eq:AB-limits}, \eqref{eq:global-sharp}, and
\eqref{eq:ratio-sharp}.  Their product tends, for fixed $m$, to
\[
  \frac{2s_m^2}{d_m}.
\]
Using \eqref{eq:s-d},
\[
  \frac{2s_m^2}{d_m}
  =
  4-2d_m.
\]
Finally $d_m\to0$, giving \eqref{eq:iterated4}.
\end{proof}

\begin{corollary}[Necessity of the inverse-square scale]
For every $c<4$ and every $\varepsilon>0$, there are indices $m,n$ in the
three-cell family such that
\[
  0<\etaq_{n,m}<\varepsilon,\qquad
  0<\dd_{n,m}^2<\varepsilon,
\]
and
\begin{equation}\label{eq:sharp-obstruction-quantified}
  \frac{\tauq_{n,m}}{\etaq_{n,m}}\,
  \dd_{n,m}^2>c.
\end{equation}
Consequently there is no universal estimate
$\tauq\le K\etaq$ with $K$ independent of the defect scale.  More generally,
a uniform estimate of the form
$\tauq\le \etaq\,g(\dd)$ cannot have
$g(\dd)=o(\dd^{-2})$ as $\dd\downarrow0$.  Thus the inverse-square power in
the scale dependence of Theorem~\ref{thm:stability} is necessary up to
multiplicative constants.
\end{corollary}

\begin{proof}
Choose $m$ sufficiently large that the fixed-$m$ scaled limit
$4-2d_m$ is larger than $c$ and that $2s_m^2$ is smaller than
$\varepsilon$.  For this fixed $m$, let $n\to\infty$.  Theorem
\ref{thm:scale-sharp} gives
\[
  \frac{\tauq_{n,m}}{\etaq_{n,m}}\dd_{n,m}^2
  \longrightarrow 4-2d_m,
  \qquad
  \dd_{n,m}^2\longrightarrow2s_m^2,
\]
while the explicit formula \eqref{eq:eta-sharp} and
$B_n^2\to0$ give $\etaq_{n,m}\to0$.  Taking $n$ sufficiently large yields
\eqref{eq:sharp-obstruction-quantified}.  The two impossibility statements
follow immediately.
\end{proof}

\section{Formal verification}

The principal results have a companion Lean 4 formalization at
\begin{center}
\url{https://github.com/Bobart0/everettian-decoherence-lean}
\end{center}
The repository is the downstream component of the pinned dependency chain
\begin{center}
\small
\nolinkurl{gleason-theorem-lean}
$\to$
\nolinkurl{quantum-foundations-lean}
$\to$\\
\nolinkurl{everettian-probability-lean}
$\to$
\nolinkurl{everettian-decoherence-lean}.
\end{center}
The audited upstream library contains formalizations of the Busch and Gleason
representation theorems, Naimark dilation, Wigner's theorem, Uhlhorn-type
uniqueness, Riedel's branch-decomposition theorem, and Kent's
contrary-inference construction.  Kent's construction is a conceptual
contrast in the upstream foundations library and is not a premise of the
present theorem.

The Lean development works with canonical coordinate Hilbert spaces
$H_n\cong\C^n$.  The paper's coordinate-free statements are invariant under
unitary change of coordinates; the repository checks the corresponding
canonical-coordinate versions.

For manuscript version 0.9, the four publication-facing landmarks are:
\begin{quote}
\small
Cut envelope:\quad
\nolinkurl{maxSubsetCommutatorOpNorm_sq_le_half_globalSq}

Exact rigidity:\quad
\nolinkurl{maxCut_saturation_iff_exactlyTwoActiveCellCommutators}

Quantitative stability:\quad
\nolinkurl{twoCellCommutatorConcentratedWithin_of_maxCut_near_saturated}

Scale obstruction:\quad
\nolinkurl{stabilitySharpnessScaleLimit_tendsto_four}
\end{quote}
The remaining exact formulas and intermediate declarations are mapped in
\nolinkurl{docs/STABILITY_RIGIDITY_SHARPNESS_MAPPING.md}; the dedicated audit
module is
\nolinkurl{EverettianDecoherence.Audit.StabilityRigiditySharpness}.
The formalization head preceding the manuscript branch is
\nolinkurl{81eb3ff53f6219749b84491513ff437fe7c545c1}; its complete CI passed,
including the library build, audit aggregate, guards, and diff-hygiene checks.
For archival submission this commit reference should be supplemented by a
tagged release and a persistent archive identifier.

\paragraph{AI-assisted formalization workflow.}
OpenAI ChatGPT was used as an interactive assistant during parts of the Lean
development, including proof-structure exploration, code drafting, and
diagnosis of compiler errors.  Every publication-facing declaration reported
here is checked by Lean itself, and the repository's continuous-integration
workflow rebuilds the library and the dedicated audit surface.  The author
reviewed the mathematical statements and proof dependencies and remains
responsible for the correctness and interpretation of the results.

\section{Discussion}

The stability theorem is dimension-free in the number of cells: its constants
do not grow with $|I|$.  The price is the explicit positive scale $\rho$.
The three-cell family shows that this scale dependence is structural rather
than merely technical.

The constants $36$ and $72$ arise from the quantitative proof architecture
and are not claimed to be optimal.  The sharpness construction identifies the
correct inverse-square scale and yields the asymptotic constant $4$ for the
explicit family, but it does not by itself determine the best universal
coefficient in the full finite-dimensional stability theorem.

The result is purely finite-dimensional and algebraic.  It requires only an
orthogonal resolution of the identity, a unitary operator, operator norms,
and finite concentration estimates; no probabilistic or dynamical structure
enters the theorem.

\section{Relation to existing operator and matrix theory}
\label{sec:related}

The classical starting point is the geometry of two projections.  Halmos'
two-subspaces theorem gives a canonical block representation for a pair of
closed subspaces, and the modern survey of B{\"o}ttcher and Spitkovsky
develops this two-projection calculus systematically
\cite{Halmos1969TwoSubspaces,BottcherSpitkovsky2010}.  Our setting differs at
the level of the basic datum: we fix a complete finite orthogonal resolution
$(P_i)_{i\in I}$ and one unitary $U$, and study simultaneously the cellwise
profile
\[
  p_i=\|[P_i,U]\|_{\op}^2
\]
and the subset functional
\[
  S\longmapsto
  \left\|
    \left[\sum_{i\in S}P_i,U\right]
  \right\|_{\op}.
\]

There is an extensive literature on commutators of two orthogonal
projections.  Li characterized the operators representable as
$[P,Q]$ for orthogonal projections, and subsequent work has studied the
corresponding representation problem and norm closures
\cite{Li2004,ShiJi2019,MarcouxRadjaviZhang2024}.  In particular, the
classical two-projection problem has the universal bound
$\|[P,Q]\|_{\op}\le1/2$; related norm identities and refinements are discussed
by Walters and Conde \cite{Conde2022}.  The occurrence of the numerical
constant $1/2$ there should not be confused with
Theorem~\ref{thm:envelope}: in our inequality the second operator is a
unitary rather than a projection, and $1/2$ multiplies the adaptive
finite-family budget
$\sum_i\|[P_i,U]\|_{\op}^2$.

Sharp commutator inequalities for arbitrary matrices provide another
methodological precedent.  The B{\"o}ttcher--Wenzel theory, for example,
bounds a Frobenius norm of a commutator and has its own extremal/equality-case
theory \cite{BottcherWenzel2008}.  Such results concern different norms and a
pair of arbitrary matrices; they do not contain the subset maximization or
the defect-profile concentration problem considered here.

The closest conceptual stability results we located for projections concern
perturbation toward an exact relation.  Walters showed that a projection
nearly orthogonal to its image under an order-two symmetry can be approximated
by one exactly orthogonal to that image, and Spitkovsky obtained an exact
distance formula in a related Hermitian-involution setting
\cite{Walters2017,Spitkovsky2018Distance}.  These results modify the
projection.  By contrast, Theorem~\ref{thm:stability} leaves both the
resolution $(P_i)$ and $U$ fixed and infers the internal shape of their
already-existing commutator profile from near equality.

This distinction also separates the present problem from the theory of
almost commuting operators.  Lin's theorem and quantitative work of Hastings
ask whether almost commuting matrices are close to exactly commuting ones
\cite{Lin1997AlmostCommuting,FriisRordam1996,Hastings2009}.  For unitaries,
topological obstructions appear
\cite{Voiculescu1983,ExelLoring1989,ExelLoring1991}, and effective or
structure-preserving versions remain active
\cite{DorOnHallKachkovskiy2025,Herrera2024,Herrera2026}.  Almost-invariant
subspaces provide a further related viewpoint
\cite{MarcouxPopovRadjavi2013}.  Again, the output in these problems is a
nearby exact relation, rather than concentration of a fixed defect profile.

A targeted literature search through these lines of work did not identify a
previous result combining the finite-family cut envelope with an exact
two-active-cell equality classification, a cardinality-free quantitative
near-equality theorem, and a three-cell inverse-square scale obstruction.
Because the envelope estimate itself is elementary once the cross-block
budget inequalities are isolated, an equivalent lemma could conceivably
exist under different block-operator or pinching terminology.  We therefore
make no absolute historical-priority claim.  The contribution asserted here
is the specific finite-family inequality together with its equality,
stability, and scale-sharpness theory.

\section{Conclusion}

For a finite orthogonal resolution and a unitary transformation, the
largest aggregate cut commutator satisfies the sharp quadratic envelope
\[
  \cmax^2\le\frac{\dd^2}{2}.
\]
Positive-defect equality is rigid: it occurs exactly when two cells, and only
two cells, carry nonzero commutator defect.  Near equality inherits a
dimension-free quantitative form.  At defect scale at least $\rho$,
near-saturation with deficit $\etaq$ forces all but
\[
  2\etaq+
  \frac{36\etaq}{(1-3\etaq)\rho^2}
\]
of the relative quadratic budget onto two cells.

An explicit three-cell family proves that the inverse-square scale cannot be
eliminated: the asymptotic scaled ratio tends to $4$.  This closes the
qualitative--quantitative--sharpness chain for the finite cut envelope and
provides a machine-checked basis for further study of near-extremizers of
finite projector--unitary commutator inequalities.

\section*{Declaration of generative AI and AI-assisted technologies in the manuscript preparation process}

During the preparation of this work, the author used OpenAI ChatGPT to assist
with organization of the manuscript, language editing, literature
organization, and consistency checks between the mathematical exposition and
the Lean formalization. After using this service, the author reviewed and
edited the content as needed and takes full responsibility for the content of
the publication.

\bibliographystyle{plain}
\bibliography{references_v03}

\end{document}
